\addcontentsline{toc}{section}{СПИСОК ИСПОЛЬЗОВАННЫХ ИСТОЧНИКОВ}

\begin{thebibliography}{99}
	
	% === Книги по 1С:Предприятие ===
	
	\bibitem{radchenko} Радченко, М. Г. 1С:Предприятие 8.3. Практическое пособие разработчика. Примеры и типовые приемы / М. Г. Радченко, Е. Ю. Хрусталева. - 3-е изд. - Москва : 1С-Паблишинг, 2023. - 983 с. - ISBN 978-5-9677-3268-3. - Текст : непосредственный.
	
	\bibitem{hrustaleva_skd} Хрусталева, Е. Ю. Разработка сложных отч-тов в 1С:Предприятии 8. Система компоновки данных / Е. Ю. Хрусталева. - 4-е изд. - Москва : 1С-Паблишинг, 2024. - 485 с. - ISBN 978-5-9677-3448-9. - Текст : непосредственный.
	
	\bibitem{hrustaleva_query} Хрусталева, Е. Ю. Язык запросов 1С:Предприятия 8.3 / Е. Ю. Хрусталева. - Москва : 1С-Паблишинг, 2025. - 378 с. - ISBN 978-5-9677-3476-2. - Текст : непосредственный.
	
	\bibitem{chrustaleva_conf} Хрусталева, Е. Ю. Расширения конфигураций. Адаптация прикладных решений с сохранением поддержки в облаках и на земле. Разработка в системе «1С:Предприятие 8.3» / Е. Ю. Хрусталева. - 3-е изд., стереотипное. - Москва : 1С-Паблишинг, 2022. - 297 с. - ISBN 978-5-9677-3171-6. - Текст : непосредственный.
	
	\bibitem{filimonova} Филимонова, Е. В. Разработка и реализация конфигураций в системе 1С:Предприятие / Е. В. Филимонова. - Москва : Синергия, 2020. - 210 с. - ISBN 978-5-4257-0502-0. - Текст : непосредственный.
	
	\bibitem{marchenko} Марченко, И. О. Разработка системы управления предприятием на платформе «1С:Предприятие 8.3» / И. О. Марченко, М. Л. Перевертайло. - Новосибирск : Изд-во НГТУ, 2018. - 116 с. - ISBN 978-5-7782-3714-8. - Текст : непосредственный.
	
	\bibitem{kashaev} Кашаев, С. М. 1С:Предприятие 8.3. Программирование и визуальная разработка / С. М. Кашаев. - Санкт-Петербург : БХВ-Петербург, 2015. - 336 с. - ISBN 978-5-9775-2806-1. - Текст : непосредственный.
	
	\bibitem{gladkiy} Гладкий, А. 1С:Бухгалтерия для начинающих / А. Гладкий. - Санкт-Петербург : Питер, 2013. - 336 с. - ISBN 978-5-496-00087-1. - Текст : непосредственный.
	
	\bibitem{gabec} Габец, А. П. Реализация прикладных задач в системе "1С:Предприятие 8.2" / А. П. Габец, Д. В. Козырев, Д. С. Кухлевский, Е. Ю. Хрусталева. - Москва : 1С-Паблишинг, 2010. - 715 с. - ISBN 978-5-9677-1387-3. - Текст : непосредственный.
	
	\bibitem{gartvitsch} Гартвич, А. В. Бухгалтерский учет в 1С:Бухгалтерия 8.3. / А. В. Гартвич. - Санкт-Петербург : БХВ-Петербург, 2015. - 285 с. - ISBN 978-5-9775-3498-7. - Текст : непосредственный.
	
	\bibitem{skorohod} Скороход, С. В. Программирование на платформе 1С:Предприятие 8.3 / С. В. Скороход. - Ростов-на-Дону ; Таганрог : Издательство Южного федерального университета, 2019. - 137 с. - ISBN 978-5-9275-3315-2. - Текст : непосредственный.
	
	\bibitem{goncharov} Гончаров, Д. И. Решение специальных прикладных задач в 1С:Предприятии 8.2 / Д. И. Гончаров, Е. Ю. Хрусталева. - Москва : 1С-Паблишинг, 2012. - 300 с. - ISBN 978-5-9677-1611-9. - Текст : непосредственный.
	
	\bibitem{kaschaev} Кашаев, С. М. Программирование в 1С:Предприятие 8.3 / С. М. Кашаев. - Санкт-Петербург : Питер, 2014. - 304 с. - ISBN 978-5-496-01234-8. - Текст : непосредственный.
	
	\bibitem{nosova} Носова, Л. С. Информационные технологии в управлении образованием / Л. С. Носова. - Челябинск : Изд-во Южно-Уральского гос. гуманитарно-пед. ун-та, 2016. - 144 с. - ISBN 978-5-906908-21-6. - Текст : непосредственный.
	
	\bibitem{titenko} Титенко, Е. А. Автоматизированные информационные системы и интеллектуальные технологии / Е. А. Титенко, Н. В. Евтеева, Л. И. Гусева, С. С. Костюков. - Курск : Юго-Западный гос. ун-т, 2013. - 131 с. - ISBN 978-5-7681-0868-7. - Текст : непосредственный.
	
	\bibitem{kochura} Кочура, Е. П. Экономическая информатика / Е. П. Кочура, В. В. Серебровский. - Курск : Юго-Западный гос. ун-т, 2013. - 161 с. - ISBN 978-5-7681-0883-0. - Текст : непосредственный.
	
	\bibitem{granina} Грянина, Е. А. Секреты профессиональной работы с «1С:Зарплата и управление персоналом 8, редакция 3». Кадровый учет, экономика и охрана труда / Е. А. Грянина, С. Г. Змиевская. - Москва : 1С-Паблишинг, 2021. - 1002 с. - ISBN 978-5-9677-3090-0. - Текст : непосредственный.
	
	\bibitem{tschistov} Чистов, Д. В. Факты хозяйственной жизни в «1С:Бухгалтерии 8» / Д. В. Чистов, В. А. Матчинов. - Изд. 2. - Москва : 1С-Паблишинг, 2025. - 477 с. - ISBN 978-5-9677-3517-2. - Текст : непосредственный.
	
	\bibitem{radchenko_chrust} Радченко, М. Г. 1С:Предприятие 8.3. Практическое пособие разработчика. Используем 1С:EDT. Издание 2, стереотипное / М. Г. Радченко, Е. Ю. Хрусталева. - Москва : 1С-Паблишинг, 2026. - 1160 с. - ISBN 978-5-9677-3623-0. - Текст : непосредственный.
	
	\bibitem{ageronok} Ажеронок, В. А. Разработка интерфейса прикладных решений на платформе «1С:Предприятие 8». Издание 2, стереотипное / В. А. Ажеронок, А. В. Островерх, М. Г. Радченко, Е. Ю. Хрусталева. - Москва : 1С-Паблишинг, 2024. - 902 с. - ISBN 978-5-9677-3359-8. - Текст : непосредственный. 
	
	\bibitem{radchenko_nachalo} Радченко, М. Г. 1С:Программирование для начинающих. Детям и родителям, менеджерам и руководителям. Разработка в системе «1С:Предприятие 8.3» / М. Г. Радченко. - 2-е изд., стереотипное. - Москва : 1С-Паблишинг, 2022. - 781 с. - ISBN 978-5-9677-3172-3. - Текст : непосредственный.
	
	\bibitem{chrustaleva} Хрусталева, Е. Ю. Технологии интеграции «1С:Предприятия 8.3» / Е. Ю. Хрусталева. - 2-е изд., стереотипное. - Москва : 1С-Паблишинг, 2023. - 503 с. - ISBN 978-5-9677-3303-1. - Текст : непосредственный.
	
	\bibitem{chrustaleva_sistema} Хрусталева, Е. Ю. Система взаимодействия. Коммуникации в бизнес-приложениях. Разработка в системе «1С:Предприятие 8.3» / Е. Ю. Хрусталева. - Москва : 1С-Паблишинг, 2019. - 130 с. - ISBN 978-5-9677-2869-3. - Текст : непосредственный.
	
	
	
	% === Управление спортивными организациями ===
	
	\bibitem{filippov} Филиппов, С. С. Менеджмент физической культуры и спорта / С. С. Филиппов. - 6-е изд., перераб. и доп. - Москва : Юрайт, 2023. - 255 с. - ISBN 978-5-534-17685-8. - Текст : непосредственный.
	
	\bibitem{matveev} Матвеев, Л. П. Теория и методика физической культуры / Л. П. Матвеев. - 5-е изд., стереотипное. - Москва : Спорт, 2025. - 520 с. - ISBN 978-5-907601-90-1. - Текст : непосредственный.
	
	\bibitem{grigorieva} Григорьева, И. И. Образование и спортивная подготовка: процессы модернизации. Вопросы и ответы. Часть 1. Организация тренировочного процесса / И. И. Григорьева, Д. Н. Черноног. - Москва : Спорт, 2016. - 336 с. - ISBN 978-5-906839-19-0. - Текст : непосредственный.
	
	\bibitem{bratkov} Братков, К. И. Менеджмент спортивных организаций / К. И. Братков, В. А. Гореликов. - Москва : Университет «Синергия», 2022. - 112 с. - ISBN 978-5-4257-0560-0. - Текст : непосредственный.
	
	\bibitem{pereverzin} Переверзин, И. И. Менеджмент спортивной организации / И. И. Переверзин. - 2-е изд., перераб. и доп. - Москва : СпортАкадемПресс, 2002. - 243 с. - ISBN 5-8134-0082-6. - Текст : непосредственный.
	
	\bibitem{zolotov} Зозуля, С. Н. Экономика физической культуры и спорта / С. Н. Зозуля, М. И. Золотов, М. М. Золотов. - Москва : Издательский центр «Академия», 2016. - 192 с. - ISBN 978-5-4468-2373-4. - Текст : непосредственный.
	
	
	\bibitem{connolly} Коннолли, Т. Базы данных : проектирование, реализация и сопровождение : теория и практика / Т. Коннолли, К. Бегг. - 3-е изд. - Москва : Вильямс, 2017. - 1439 с. - ISBN 978-5-8459-2020-1. - Текст : непосредственный.
	
	
	% === Программная инженерия и проектирование ===
	
	
	
	\bibitem{martin_clean} Мартин, Р. Чистый код: создание, анализ и рефакторинг / Р. Мартин. - Санкт-Петербург : Питер, 2019. - 464 с. - ISBN 978-5-4461-0960-9. - Текст : непосредственный.
	
	\bibitem{martin_clean_arch} Мартин, Р. Чистая архитектура : искусство разработки программного обеспечения / Р. Мартин. - Санкт-Петербург : Питер, 2018. - 401 с. - ISBN 978-5-4461-0772-8. - Текст : непосредственный.
	
	\bibitem{fowler} Фаулер, М. Рефакторинг. Улучшение проекта существующего кода / М. Фаулер. - Москва ; Санкт-Петербург : Диалектика, 2019. - 442 с. - ISBN 978-5-9909445-1-0. - Текст : непосредственный.
	
	\bibitem{booch} Буч, Г. Язык UML. Руководство пользователя / Г. Буч, Дж. Рамбо, И. Якобсон. - 2-е изд. - Москва : ДМК-Пресс, 2015. - 493 с. - ISBN 5-94074-334-X. - Текст : непосредственный.
	
	\bibitem{belik} Белик, А. Г. Проектирование и архитектура программных средств / А. Г. Белик, В. Н. Цыганенко. - Омск : Изд-во ОмГТУ, 2016. - 94 с. - ISBN 978-5-8149-2258-8. - Текст : непосредственный.
	
	
	
	% === Проектирование интерфейсов ===
	
	\bibitem{krug} Круг, С. Не заставляйте меня думать / С. Круг. - 3-е изд. - Москва : Эксмо, 2025. - 256 с. - ISBN 978-5-699-91492-0. - Текст : непосредственный..
	
	\bibitem{norman} Норман, Д. А. Дизайн привычных вещей / Д. А. Норман. - Москва : Манн, Иванов и Фербер, 2026. - 384 с. - ISBN 978-5-00195-363-0. - Текст : непосредственный.
	
	\bibitem{cooper} Купер, А. Интерфейс. Основы проектирования взаимодействия / А. Купер. - 4-е изд. - Санкт-Петербург : Питер, 2024. - 720 с. - ISBN 978-5-4461-0877-0. - Текст : непосредственный.
	
	% === Тестирование ПО ===
	
	\bibitem{reschka} Рэшка, Д. Тестирование программного обеспечения. Внедрение, управление и эксплуатация / Д. Рэшка, Э. Дастин, Дж. Пол. - Москва : Лори, 2012. - 567 с. - ISBN 978-5-85582-318-9. - Текст : непосредственный.
	
	\bibitem{myers} Майерс, Г. Искусство тестирования программ / Г. Майерс, Т. Баджетт, К. Сандлер. - 3-е изд. - Санкт-Петербург : Диалектика, 2020. - 271 с. - ISBN 978-5-907203-66-2. - Текст : непосредственный.
	
	\bibitem{keyner} Кейнер, К. Тестирование программного обеспечения : контекстно ориентированный подход / К. Кейнер, Д. Бах, Б. Петтикорд. - Санкт-Петербург : Питер, 2025. - 352 с. - ISBN 978-5-4461-2165-6. - Текст : непосредственный.
	
	\bibitem{gagloeva} Гаглоева, И. Э. Методы системного анализа, принятия решений и обработки информации с помощью компьютерных программ / И. Э. Гаглоева, Ю. В. Саханский, М. А. Ковалева, М. В. Волик. - Москва : Прометей, 2024. - 90 с. - ISBN 978-5-00172-615-9. - Текст : непосредственный.
	
	% === Документооборот и автоматизация ===
	
	\bibitem{rogogin} Рогожин, М. Ю. Настольная книга ответственного за делопроизводство / М. Ю. Рогожин. - Москва : Проспект, 2026. - 127 с. - ISBN 978-5-392-45020-6. - Текст : непосредственный.
	
	\bibitem{kuznetsov_doc} Кузнецов, И. Н. Документационное обеспечение управления. Документооборот и делопроизводство / И. Н. Кузнецов. - 5-е изд., перераб. и доп. - Москва : Юрайт, 2025. - 425 с. - ISBN 978-5-534-20025-6. - Текст : непосредственный.
	
	% === Государственные стандарты ===
	
	\bibitem{gost_19_701} ГОСТ 19.701-90. Единая система программной документации. Схемы алгоритмов, программ, данных и систем. - Москва : Стандартинформ, 2010. Текст : непосредственный.
	
	
	\bibitem{gost_34_602} ГОСТ 34.602-2020. Информационные технологии. Комплекс стандартов на автоматизированные системы. Техническое задание на создание автоматизированной системы. - Москва : Стандартинформ, 2020. Текст : непосредственный.
	
	
	\bibitem{gost_34_201} ГОСТ 34.201-2020. Информационные технологии. Комплекс стандартов на автоматизированные системы. Виды, комплектность и обозначение документов при создании автоматизированных систем. - Москва : Стандартинформ, 2020. Текст : непосредственный.
	
	
	
	
	
	
	% === Электронные ресурсы ===
	
	\bibitem{1c_platform} 1С:Предприятие 8.3. Документация по платформе : официальный сайт. - URL: https://v8.1c.ru/platforma/ (дата обращения: 06.05.2026).
	
	\bibitem{1c_its} Информационно-технологическое сопровождение 1С : официальный сайт. - URL: https://its.1c.ru/ (дата обращения: 06.05.2026).
	
	\bibitem{vanessa} Vanessa Automation : официальный сайт. - URL: https://github.com/Pr-Mex/vanessa-automation (дата обращения: 06.05.2026).
	
	\bibitem{plantuml} PlantUML : официальный сайт. - URL: https://plantuml.com/ru/ (дата обращения: 08.05.2026).
	
	\bibitem{drawio} Draw.io : официальный сайт. - URL: https://www.diagrams.net/ (дата обращения: 08.05.2026).
	
	\bibitem{minobr} Министерство науки и высшего образования РФ : официальный сайт. - URL: https://fgos.ru/ (дата обращения: 08.05.2026).
	
	\bibitem{minsport} Министерство спорта Российской Федерации : официальный сайт. - URL: https://www.minsport.gov.ru/ (дата обращения: 04.05.2026).
	
	\bibitem{evsk} Единая всероссийская спортивная классификация : официальный сайт. - URL: https://www.minsport.gov.ru/sport/podgotovka/edinaya-vserossiyskaya-sportivnaya-klassifikaciya/ (дата обращения: 04.05.2026).
	
	
	
\end{thebibliography}